\documentclass[12pt]{article}
\usepackage[utf8]{inputenc}
\usepackage{amsmath}
\usepackage{graphicx}
\usepackage{natbib}
\usepackage{hyperref}
\usepackage{lipsum} % For dummy text, you can remove this in your final document

\title{Nanomaterials in Biomedical Applications: Drug Delivery Systems, Diagnostics, and Imaging Technologies}
\author{Your Name}
\date{\today}

\begin{document}

\maketitle

\begin{abstract}
Nanomaterials have emerged as versatile tools in biomedical applications, offering significant advancements in drug delivery, diagnostics, and imaging technologies. This paper reviews the types of nanomaterials used, discusses their applications in biomedicine, examines the benefits, and analyzes current challenges and limitations.
\end{abstract}

\section{Introduction}
Nanotechnology has revolutionized biomedical research by providing innovative solutions for drug delivery, disease diagnostics, and medical imaging. Nanomaterials, due to their unique physical and chemical properties at the nanoscale, offer promising opportunities in improving therapeutic efficacy, enhancing diagnostic sensitivity, and enabling precise imaging techniques.

\section{Nanomaterial Types}
\subsection{Organic Nanomaterials}
Organic nanoparticles such as liposomes, dendrimers, and polymer nanoparticles are widely used for drug encapsulation, targeted delivery, and controlled release systems \citep{sun2014nanotechnology, alexis2008nanomedicine}.

\subsection{Inorganic Nanomaterials}
Inorganic nanoparticles including quantum dots, gold nanoparticles, and magnetic nanoparticles serve as contrast agents for imaging, carriers for therapeutic agents, and platforms for biosensing applications \citep{medintz2008quantum, giljohann2010gold}.

\section{Biomedical Applications}
\subsection{Drug Delivery Systems}
Nanomaterial-based drug delivery systems enable targeted delivery, prolonged circulation time, and controlled release of therapeutic agents, minimizing side effects and improving patient compliance \citep{wang2018nanoparticle, allen2017nanoparticle}.

\subsection{Diagnostics}
Nanotechnology enhances diagnostic capabilities through improved sensitivity and specificity in detecting biomarkers, pathogens, and disease conditions, facilitating early disease detection and personalized medicine approaches \citep{cao2019nanomaterials, kwizera2020nanomaterials}.

\subsection{Imaging Technologies}
Nanomaterials serve as contrast agents in various imaging modalities such as magnetic resonance imaging (MRI), computed tomography (CT), and fluorescence imaging, offering enhanced resolution and multifunctional imaging capabilities \citep{choi2008gold, tsourkas2005magneto}.

\section{Benefits}
\subsection{Therapeutic Advantages}
Nanomaterials improve drug solubility, stability, and bioavailability, allowing for targeted therapy and reduced systemic toxicity \citep{dreaden2012nanotechnology, moore2014nanotechnology}.

\subsection{Diagnostic Precision}
Enhanced sensitivity of nanomaterial-based biosensors and imaging agents enables early disease detection and real-time monitoring of therapeutic responses \citep{bai2017nanomaterials, medintz2005nanomaterials}.

\section{Challenges}
\subsection{Biocompatibility and Safety}
Concerns over potential toxicity, immunogenicity, and long-term effects of nanomaterials in biological systems \citep{chen2013review, pelaz2012biocompatibility}.

\subsection{Regulatory Hurdles}
Navigating regulatory pathways for clinical translation of nanomedicine products, ensuring safety, efficacy, and quality standards \citep{hirsch2013regulation, saffari2019regulatory}.

\section*{References}
\bibliographystyle{plainnat}
\bibliography{prompt11_35}

\end{document}
